\documentclass[10pt,a4paper]{amsart}
\usepackage[margin=19mm]{geometry}
\usepackage{amsmath,amssymb,mathtools}
\usepackage[scr=rsfs,cal=euler]{mathalfa}
\usepackage{xfrac}
\usepackage[unicode,hidelinks,bookmarks=false]{hyperref}
\newcommand{\ie}{i.e.\ }
\newcommand{\cf}{cf.\ }
\newcommand{\resp}{respectively\ }
\newcommand{\Subsection}{\subsection}
\providecommand{\nobreakdash}{\nobreak\hbox{-}}
\setlength{\emergencystretch}{2em}
\multlinegap0pt
\usepackage{bm}
\usepackage{bbm}
\usepackage{dsfont}
\usepackage{xspace}
\usepackage{cancel}
\usepackage{mathtools}
\usepackage{amsthm}
\newtheorem{lemma}{Lemma}
\title[Risk-Sensitive Learning in Population Games under Extreme Ev]{Risk-Sensitive Learning in Population Games under Extreme Events: Bifurcations and Chaotic Dynamics}
\author{Konstantinos Metaxas, Themistoklis P. Sapsis}
\date{Source: arXiv:2606.29967v2 (CC BY 4.0)}
\begin{document}
\maketitle

\noindent\emph{Source and licence.} Risk-Sensitive Learning in Population Games under Extreme Events: Bifurcations and Chaotic Dynamics. Version arXiv:2606.29967v2 (CC BY 4.0), distributed under the Creative Commons Attribution 4.0 International licence, as recorded in the arXiv licence metadata for this exact version and shown on the abstract page https://arxiv.org/abs/2606.29967v2. Full rights notice: licenses/arxiv-2606.29967-CC-BY-4.0.txt.\par\smallskip

\noindent\textit{Editorial scope.} Counted unit: the Lemma whose source label is \texttt{lem\_delta} in the article (source0.tex lines 977--979, source label \texttt{lem\_delta}), delivered together with the article's own complete original proof of it (source0.tex lines 981--1016, one \texttt{proof} environment of 36 lines). No mathematics is written, repaired or re-derived: statement and proof are the article's own text line for line. Required context retained uncounted: sys (lines 202--208). Every definition, display and label of those lines is retained unchanged. Omitted from this package and not redistributed: 1--201, 209--976, 980--980, 1017--1216. Nothing inside a retained excerpt is omitted or altered: every retained body is delivered exactly as the article's lines are written, so no omission receipt of any kind accompanies this package. Citations: the retained excerpts carry no citation command at all, so no bibliography entry of the article is transcribed and no reference is redistributed. Reference adaptation: none was needed and none was made; every reference inside the delivered unit resolves inside it, so no unresolved reference or citation is produced. Printed slips kept literal: no printed word, symbol, hypothesis or number of the counted statement or of its proof was altered. Layout and class: the article's own class is \texttt{article} with packages including multirow, inputenc, fontenc, geometry, tikz, soul, enumitem, subfigure, array, float, booktabs, color, amsmath, amssymb; the wrapper is the lane's pinned \texttt{amsart} editorial wrapper at 10pt on A4, which restates the theorem-like environments the retained text needs and adds the article's own macro definitions that the retained text needs (none), with extra wrapper packages (bm, bbm, dsfont, xspace, cancel, mathtools). The delivered numbering is the unit's own. Assets: no retained span contains a figure environment, a graphics-inclusion command, a PDF-page inclusion, a table or a TikZ picture; no image file is copied into this package, the package declares no asset dependency and no image file is redistributed with it. Licence: version arXiv:2606.29967v2 of this article is distributed under the Creative Commons Attribution 4.0 International licence, as recorded in the arXiv licence metadata for this exact version and shown on the abstract page https://arxiv.org/abs/2606.29967v2. A full rights notice is delivered at licenses/arxiv-2606.29967-CC-BY-4.0.txt. Characters: every retained excerpt is pure ASCII, so the delivered engine, XeLaTeX with its default Latin Modern text font, reports no missing character.
\begin{align}
x_{n+1}&=\frac{x_n}{x_n+(1-x_n)\exp\!\left(
a\left(x_n-b+cq_n \right)
\right)} \nonumber\\
q_{n+1}&=(1-k)q_n+k(p_0+p_1x_n).
\label{sys}
\end{align}

\begin{lemma} \label{lem_delta}
If $x_0\in(0,1)$ and $q_0\in[0,1],$ then there exists a $\delta_1\in(0,1)$ such that $x_n<1-\delta_1$ for all $n\geq 0.$ Furthermore, if $b>cp_0,$ then there is a $\delta\in(0,1)$ such that $x_n\in(\delta,1-\delta)$ for all $n.$
\end{lemma}

\begin{proof}
For this proof it is more convenient to work with the logits $y_n=\log \left(x_n/(1-x_n) \right).$ Then, using \eqref{sys} we have
\[
y_{n+1}=y_n-a(x_n-b+cq_n).
\]
Let $\varepsilon_1=(1-b)/2.$ If \(x_n\geq 1-\varepsilon_1\), then $x_n-b+cq_n>0$ and, thus, $y_{n+1}<y_n$, implying $x_{n+1}<x_n.$ By continuity of the \(x\)-component of the map on the compact set \([0,1-\varepsilon_1]\times[0,1]\), there exists \(\mu<1\) such that  $x_n\leq \mu<1$ whenever $(x_{n-1},q_{n-1})\in[0,1-\varepsilon_1]\times[0,1]$. As a result, there is a $1>\delta_1>0$ such that $x_n<1-\delta_1$ for all $n.$

Now suppose that $b>cp_0.$ Let $\bar q\in(p_0,b/c)$ and \(\varepsilon_2>0\) small enough so that 
\[
p_0+p_1\varepsilon_2<\bar q, \; \; \varepsilon_2 - b + c\bar q<0.
\]
Let $\theta=p_0+p_1\varepsilon_2<\bar q$ and $\beta=b-c\bar q-\varepsilon_2>0$. If \(x_n\leq\varepsilon_2\) and \(q_n\leq\bar q\), then
\[
x_n-b+cq_n\leq \varepsilon_2 - b+c\bar q=-\beta,
\]
and hence $y_{n+1}\geq y_n+a\beta.$ Thus, once \(x_n\) is sufficiently small and \(q_n\leq\bar q\), the \(x\)-coordinate is pushed away from zero. It remains to control the indices for which \(x_n\leq\varepsilon_2\) but \(q_n\) may still be larger than \(\bar q\). Consider $N$ indices such that \(x_s,x_{s+1},\ldots,x_{s+N-1}\leq\varepsilon_2\). Then 
\[
q_{s+N}
\leq
(1-k)^N q_s+\bigl(1-(1-k)^N\bigr)\theta
\leq
(1-k)^N+\bigl(1-(1-k)^N\bigr)\theta.
\]
But since \(\theta<\bar q\), there is an $L\geq 1$ such that 
\[
(1-k)^L+\bigl(1-(1-k)^L\bigr)\theta\leq \bar q.
\]
As a result, if $L\leq N-1$, then \(q_{s+L}\leq\bar q\). Moreover, $q_{s+L+1}< \bar q$ and, thus, $q_{s+j}\leq \bar q$ for every $N-1\geq j\geq L,$ since $x_{s+j}\leq \varepsilon_2$ continues to hold. That is, after the first $L$ steps, $y_n,$ and hence $x_n$, begins to increase, as long as $x_n\leq \varepsilon_2$, suggesting that there cannot be infinitely many consecutive indices with $x_n<\varepsilon_2$.

Now let $A=a(1-b+c)$ and $y_{\varepsilon}=\log \left( \varepsilon_2/( 1-\varepsilon_2) \right)$. Then, $y_{n+1}\geq y_n-A$ for all $n$. Consider a maximal block of consecutive indices starting at $s>0$ and of length $N$ during which \(x_n\leq\varepsilon_2\). First, we have $x_{s-1}>\varepsilon_2$, implying $y_s\geq y_{\varepsilon}-A$. Based on the previous discussion, the logits $y_n$ can decrease by at most $LA,$ as after the first $L$ steps, $q_n \leq \bar q$ and, so, $y_n$ increases. Therefore, within any such block we have
\[
y_{n} \geq y_\varepsilon-(L+1)A,\; n=s,s+1,\dots,s+N-1.
\]
If there is a block with $s=0,$ then one has $y_n\geq y_0-LA.$ Outside the region \(x\leq\varepsilon_2\), we have \(y>y_{\varepsilon}\) and, thus, there exists a constant $Y$ such that $y_n\geq Y$ for all $n.$ Equivalently, there exists $\delta_2 \in(0,1)$ such that $x_n > \delta_2$ for all $n$. Combining this result with the upper bound, we obtain the claim.

\end{proof}

\end{document}
